\section*{The Propers: Themes and Movement}

\begin{massbox}{From effective word to rooted fruit}
The Lectionary's correlated strand moves from God's effective word through watered earth to the sower and the varied reception of seed. Romans supplies a distinct apostolic horizon in which even those with firstfruits still groan while creation awaits liberation, while the offering, Communion, and final prayers ask that received gifts sanctify a people whose fruit matures through repeated participation.
\end{massbox}

\begin{studybox}{Light, word, groaning, gift}
\textbf{1. Light seeks wanderers --- Entrance; Collect.} The just speaker hopes to behold, and truth calls Christians back to their name.\par\smallskip
\textbf{2. Word descends and is heard --- First Reading; Psalm; Acclamation; Gospel.} Rain, harvest, and seed disclose gift and differentiated reception; the long and short forms do not proclaim identical material.\par\smallskip
\textbf{3. Creation groans in hope --- Second Reading.} Paul's semi-continuous strand joins present suffering, firstfruits, and bodily redemption.\par\smallskip
\textbf{4. Gift becomes sanctification --- Prayer over Offerings; Communion; Prayer after Communion.} The appointed prayers and antiphons nourish enduring fruit.
\end{studybox}

\subsection*{1. Truth illuminates before the soil answers\enspace\elementref{Entrance; Collect}}

The Entrance's hope to behold divine glory and the Collect's light of truth begin with reception. The wandering are shown a path before being asked to walk it. The prayer then makes Christian identity morally serious: certain actions contradict the name, while others accord with it. This Week XV movement prepares hearers for the Gospel's differentiated response, though the oration serves every cycle rather than Matthew 13 alone.

\subsection*{2. The correlated strand: descent, watering, sowing, fruit\enspace\elementref{First Reading; Psalm; Acclamation; Gospel}}

Isaiah's word is effective because it is sent by God, not because the earth controls the sky. Psalm 65 answers with land visited, watered, softened, and crowned; its Luke-marked response names fruitful reception. The composed acclamation gives the Sunday's explicit Christological identification: Christ is the sower of the word. Matthew then refuses a simplistic success story by displaying loss, shallowness, choking pressure, and abundant fruit.

\subsection*{3. Firstfruits do not abolish groaning\enspace\elementref{Second Reading}}

Romans belongs to its own sequence. Paul does not say that present suffering is unreal or that believers already possess bodily consummation. Creation waits, labors, and hopes; Spirit-bearing Christians groan too. Chrysostom's direct exposition preserves the passage's personification and its movement toward incorruption.

\subsection*{4. Repetition serves rooted transformation\enspace\elementref{Prayer over Offerings; Communion; Prayer after Communion}}

The Prayer over the Offerings asks that ecclesial gifts advance sanctification; Communion A names dwelling and praise, while Communion B names Johannine abiding. The final prayer asks that repeated participation deepen saving effect. Neither antiphon makes sacramental fruit automatic; each gives a distinct textual path into the same concluding petition.
